% Related work
\section{Related Work}
\label{sec:related}

Concept erasure methods differ mainly in the function class against which they erase and in the geometry they assume. Linear methods such as INLP \citep{ravfogel-etal-2020-null} and LEACE \citep{belrose-etal-2023-leace} target linearly decodable concept information. INLP removes successively exposed linear directions by iterative nullspace projection, while LEACE gives the strongest closed-form linear baseline in our comparisons. These methods are essential reference points, but they do not directly address the case where a fresh nonlinear probe can still recover the concept after linear erasure.

That recoverability concern is central to our framing. Elazar and Goldberg \citep{elazar-goldberg-2018-adversarial} showed that apparent success against one learned remover or classifier does not imply that the protected information is gone for a newly trained downstream model. Amnesic probing \citep{elazar-etal-2021-amnesic} and related intervention-based analyses adopt the same conservative lesson: the meaningful question is what remains recoverable after the intervention, not only what the original scorer no longer captures.

Two nonlinear erasure methods anchor our comparisons. IGBP \citep{iskander-etal-2023-shielded} is the closest prior baseline in mechanism: it iteratively trains a nonlinear probe and retracts representations toward the current decision boundary. Obliviator \citep{akbari-etal-2025-obliviator} is the strongest published baseline in headline numbers: it formulates erasure in RKHS with a Hilbert--Schmidt independence criterion (HSIC) penalty between the representation and the target attribute, and gradually morphs the feature space via random Fourier feature eigenvalue projection (RFF-EVP). Both methods are supervised on the target concept and unsupervised on control concepts (the same regime as our method); they differ from \TanGCE{} in how the update is constructed, not in what supervision they use. We report INLP, IGBP, and Obliviator under the same trajectory-probe protocol as every other method in Tab.~\ref{tab:main_results_unified}. Our paper differs in two ways. First, \AmbGCE{} is introduced only as the direct ambient-space ablation that keeps the iterative nonlinear scorer loop but removes the manifold constraint. Second, \TanGCE{} is the paper's main method: it projects the erasure update onto a locally estimated tangent space, making the geometric hypothesis itself the main point of comparison.
